%% kapa: exam=2 round=33 year=2022 period=2 subject=appraisal_theory subject_ko=감정평가이론 seq=1 number=1 points=40
%% source: Q-Net 감정평가사 기출문제 (cst00309) artlSeq=5210838 file=raw/1008/5210838/2022년 제33회 감정평가사 2차시험 2교시 문제지(감정평가이론).pdf
%% conversion: kordoc 4.24.51 -> tools/gen2.py (auto); review=proofread-gpt-astra; flags=none
\begin{longquestion}{【문제 1】}
최근 지식재산권에 대한 관심이 높아지면서 지식재산권에 대한 감정평가 수요도 증가하고 있다. 지식재산권 감정평가와 관련하여 다음 물음에 답하시오. (40점) \par
\medskip\noindent 물음 1) 감정평가 실무기준상 지식재산권의 개념 및 종류, 가격자료에 대해 설명하시오. (10점) \par
\medskip\noindent 물음 2) 감정평가 3방식의 성립 근거와 각 방식 간의 관계에 대해 설명하시오. (10점) \par
\medskip\noindent 물음 3) 감정평가 실무기준상 감정평가 3방식에 따른 지식재산권의 평가방법을 설명하고, 각 방식 적용 시 유의사항에 대해 설명하시오. (20점) \par
\end{longquestion}
